% KALUSAGAP v37, commit 9fa1cf359339dc92fa26de2d189554771391229d
% Add to the manuscript preamble: \usepackage{longtable,booktabs,amssymb,array}
% Paste or \input this file in place of the original table, outside any table float.
% The manuscript counter supplies the table number, e.g. Table 4.1.
\begingroup
\small
\setlength{\tabcolsep}{2pt}
\renewcommand{\arraystretch}{1.16}
\newcommand{\KalGrant}{\ensuremath{\checkmark}}
\begin{longtable}{@{}>{\raggedright\arraybackslash}p{0.38\linewidth}>{\centering\arraybackslash}p{0.085\linewidth}*{5}{>{\centering\arraybackslash}p{\dimexpr(0.425\linewidth-14\tabcolsep)/5\relax}}>{\centering\arraybackslash}p{0.11\linewidth}@{}}
\caption{User Role and Functionality Matrix}\label{tab:user-role-functionality-matrix}\\
\toprule
\textbf{Permission} & {\scriptsize\textbf{Admin}} & {\scriptsize\textbf{MHO}} & {\scriptsize\textbf{PHN}} & {\scriptsize\textbf{HS}} & {\scriptsize\textbf{RHU}} & {\scriptsize\textbf{BHW}} & {\scriptsize\textbf{Resident}} \\
\midrule
\endfirsthead
\multicolumn{8}{@{}l}{\tablename~\thetable\ (continued)}\\
\toprule
\textbf{Permission} & {\scriptsize\textbf{Admin}} & {\scriptsize\textbf{MHO}} & {\scriptsize\textbf{PHN}} & {\scriptsize\textbf{HS}} & {\scriptsize\textbf{RHU}} & {\scriptsize\textbf{BHW}} & {\scriptsize\textbf{Resident}} \\
\midrule
\endhead
\endfoot
\bottomrule
\endlastfoot
\multicolumn{8}{@{}l}{\textbf{Resident Management}}\\*
% residents.directory.view
View resident directory {\footnotesize(View)} & \KalGrant & \KalGrant & \KalGrant & \KalGrant & \KalGrant & -- & -- \\
% residents.profile.view
View resident profile {\footnotesize(View)} & \KalGrant & \KalGrant & \KalGrant & \KalGrant & \KalGrant & \KalGrant & \KalGrant \\
% residents.create
Add resident {\footnotesize(Create)} & \KalGrant & \KalGrant & \KalGrant & \KalGrant & \KalGrant & \KalGrant & -- \\
% residents.edit
Edit resident information {\footnotesize(Edit)} & \KalGrant & -- & \KalGrant & \KalGrant & \KalGrant & \KalGrant & -- \\
% residents.registration.approve
Approve resident registration {\footnotesize(Approve)} & \KalGrant & -- & \KalGrant & \KalGrant & -- & -- & -- \\
% residents.transfer.approve
Approve transfer of residency {\footnotesize(Approve)} & \KalGrant & \KalGrant & \KalGrant & \KalGrant & -- & -- & -- \\
% residents.archive
Archive/deactivate resident record {\footnotesize(Delete)} & \KalGrant & -- & -- & -- & -- & -- & -- \\
\midrule
\multicolumn{8}{@{}l}{\textbf{Household Profiling}}\\*
% households.view
View household profiles {\footnotesize(View)} & \KalGrant & -- & \KalGrant & \KalGrant & -- & \KalGrant & -- \\
% households.create
Create household profile {\footnotesize(Create)} & \KalGrant & -- & \KalGrant & -- & -- & \KalGrant & -- \\
% households.verify
Verify household information {\footnotesize(Approve)} & \KalGrant & -- & -- & \KalGrant & -- & -- & -- \\
\midrule
\multicolumn{8}{@{}l}{\textbf{Consultation}}\\*
% consultation.requests.view
View consultation requests {\footnotesize(View)} & \KalGrant & -- & \KalGrant & \KalGrant & \KalGrant & -- & \KalGrant \\
% consultation.requests.approve
Approve consultation requests {\footnotesize(Approve)} & \KalGrant & -- & -- & -- & -- & -- & -- \\
% consultation.conduct
Conduct consultation {\footnotesize(Create)} & \KalGrant & -- & \KalGrant & \KalGrant & \KalGrant & -- & -- \\
% consultation.findings.record
Record consultation findings {\footnotesize(Create)} & \KalGrant & -- & \KalGrant & \KalGrant & \KalGrant & -- & -- \\
% consultation.records.update
Update consultation records {\footnotesize(Edit)} & \KalGrant & -- & \KalGrant & \KalGrant & -- & -- & -- \\
% consultation.history.view
View consultation history {\footnotesize(View)} & \KalGrant & \KalGrant & \KalGrant & \KalGrant & \KalGrant & -- & \KalGrant \\
\midrule
\multicolumn{8}{@{}l}{\textbf{Medical Certificates}}\\*
% certificates.request
Request medical certificates {\footnotesize(Create)} & \KalGrant & \KalGrant & \KalGrant & -- & \KalGrant & -- & -- \\
\midrule
\multicolumn{8}{@{}l}{\textbf{Triage}}\\*
% triage.view
View triage records {\footnotesize(View)} & \KalGrant & -- & \KalGrant & \KalGrant & \KalGrant & -- & -- \\
% triage.perform
Perform triage {\footnotesize(Create)} & \KalGrant & -- & -- & \KalGrant & \KalGrant & -- & -- \\
% triage.assessment.update
Update triage assessment {\footnotesize(Edit)} & \KalGrant & -- & -- & \KalGrant & \KalGrant & -- & -- \\
% triage.forward
Forward patient to Health Supervisor {\footnotesize(Approve)} & \KalGrant & -- & -- & \KalGrant & \KalGrant & -- & -- \\
\midrule
\multicolumn{8}{@{}l}{\textbf{Referrals}}\\*
% referrals.view
View referrals {\footnotesize(View)} & \KalGrant & \KalGrant & \KalGrant & \KalGrant & \KalGrant & -- & \KalGrant \\
% referrals.create
Create referral {\footnotesize(Create)} & \KalGrant & -- & \KalGrant & \KalGrant & -- & -- & -- \\
% referrals.approve
Approve referral {\footnotesize(Approve)} & \KalGrant & -- & \KalGrant & \KalGrant & \KalGrant & -- & -- \\
% referrals.reject
Reject referral {\footnotesize(Approve)} & \KalGrant & -- & \KalGrant & \KalGrant & \KalGrant & -- & -- \\
% referrals.assign
Assign referral {\footnotesize(Edit)} & \KalGrant & -- & \KalGrant & \KalGrant & -- & -- & -- \\
% referrals.status.update
Update referral status {\footnotesize(Edit)} & \KalGrant & -- & \KalGrant & \KalGrant & \KalGrant & -- & -- \\
% referrals.history.view
View referral history {\footnotesize(View)} & \KalGrant & \KalGrant & \KalGrant & \KalGrant & \KalGrant & -- & \KalGrant \\
\midrule
\multicolumn{8}{@{}l}{\textbf{Follow-up Management}}\\*
% followups.view
View follow-up schedules {\footnotesize(View)} & \KalGrant & \KalGrant & \KalGrant & \KalGrant & -- & -- & \KalGrant \\
% followups.create
Create follow-up schedule {\footnotesize(Create)} & \KalGrant & -- & \KalGrant & \KalGrant & -- & -- & -- \\
% followups.edit
Edit follow-up schedule {\footnotesize(Edit)} & \KalGrant & -- & \KalGrant & \KalGrant & -- & -- & -- \\
% followups.complete
Mark follow-up as completed {\footnotesize(Edit)} & \KalGrant & -- & \KalGrant & \KalGrant & -- & -- & -- \\
% followups.history.view
View follow-up history {\footnotesize(View)} & \KalGrant & \KalGrant & \KalGrant & \KalGrant & -- & -- & \KalGrant \\
% followups.notify
Send follow-up notification {\footnotesize(Create)} & \KalGrant & -- & \KalGrant & \KalGrant & -- & -- & -- \\
\midrule
\multicolumn{8}{@{}l}{\textbf{Health Services}}\\*
% services.view
View health services {\footnotesize(View)} & \KalGrant & \KalGrant & \KalGrant & \KalGrant & \KalGrant & \KalGrant & -- \\
% services.create
Add health service {\footnotesize(Create)} & \KalGrant & \KalGrant & \KalGrant & \KalGrant & -- & -- & -- \\
% services.edit
Edit health service {\footnotesize(Edit)} & \KalGrant & \KalGrant & \KalGrant & \KalGrant & -- & -- & -- \\
% services.delete
Delete health service {\footnotesize(Delete)} & \KalGrant & \KalGrant & \KalGrant & \KalGrant & -- & -- & -- \\
\midrule
\multicolumn{8}{@{}l}{\textbf{Reports}}\\*
% reports.view
View reports {\footnotesize(View)} & \KalGrant & \KalGrant & \KalGrant & \KalGrant & -- & -- & -- \\
% reports.generate
Generate reports {\footnotesize(Create)} & \KalGrant & \KalGrant & \KalGrant & \KalGrant & -- & \KalGrant & -- \\
% reports.export
Export reports {\footnotesize(Create)} & \KalGrant & \KalGrant & \KalGrant & \KalGrant & -- & -- & -- \\
% reports.analytics.view
View analytics {\footnotesize(View)} & \KalGrant & \KalGrant & \KalGrant & \KalGrant & -- & -- & -- \\
\midrule
\multicolumn{8}{@{}l}{\textbf{User / Account Management}}\\*
% accounts.view
View user accounts {\footnotesize(View)} & \KalGrant & -- & -- & -- & -- & -- & -- \\
% accounts.create
Create user accounts {\footnotesize(Create)} & \KalGrant & \KalGrant & \KalGrant & \KalGrant & -- & -- & -- \\
% accounts.personnel.approve
Approve health personnel accounts {\footnotesize(Approve)} & \KalGrant & -- & \KalGrant & \KalGrant & -- & -- & -- \\
% accounts.edit
Edit user accounts {\footnotesize(Edit)} & \KalGrant & -- & -- & -- & -- & -- & -- \\
% accounts.deactivate
Deactivate user accounts {\footnotesize(Delete)} & \KalGrant & -- & -- & -- & -- & -- & -- \\
% accounts.delete
Delete user accounts {\footnotesize(Delete)} & \KalGrant & -- & -- & -- & -- & -- & -- \\
% accounts.access.reset
Reset user access {\footnotesize(Edit)} & \KalGrant & -- & -- & -- & -- & -- & -- \\
% accounts.roles.manage
Manage roles and permissions {\footnotesize(Approve)} & \KalGrant & -- & -- & -- & -- & -- & -- \\
\midrule
\multicolumn{8}{@{}l}{\textbf{System Management}}\\*
% system.activity.view
View system activity logs {\footnotesize(View)} & \KalGrant & -- & -- & -- & -- & -- & -- \\
% system.audit.view
View audit logs {\footnotesize(View)} & \KalGrant & -- & -- & -- & -- & -- & -- \\
% system.settings.manage
Manage system settings {\footnotesize(Edit)} & \KalGrant & -- & -- & -- & -- & -- & -- \\
% system.notifications.manage
Manage notification settings {\footnotesize(Edit)} & \KalGrant & -- & -- & -- & -- & -- & -- \\
\end{longtable}
\raggedright
\noindent\textit{Legend.} \KalGrant{} denotes a default grant in the Administrator's Roles and Permissions catalogue. A dash denotes a permission off by default. Parentheses show the action column used in that screen. MHO means Municipal Health Office, PHN means Public Health Nurse, HS means Health Supervisor, RHU denotes Rural Health Unit Personnel, and BHW means Barangay Health Worker.

\par\medskip
\noindent\textit{Authorization and scope in KALUSAGAP v37.} The screen overlays saved database values on the catalogue defaults. The API applies its own feature-role checks first. On routes that check persisted grants, a role with no saved entries falls back to feature-role authorization. Once entries exist for the role, a missing required grant can deny the API request even if the UI still shows a catalogue default. Service and database rules also enforce record ownership, geography, workflow state, and staff assignments. The table describes the source defaults. For example, health-service catalogue reads use a feature-role gate without testing the View health services toggle. Certificate decisions also use a feature-role gate separate from the Request medical certificates toggle.

\begin{description}
\item[Admin.] Admin can manage account provisioning, role and assignment changes, lifecycle decisions, access recovery, permanent account deletion, and role-permission settings. The API and database preserve the administrator's permission-management grant and audit changes to the matrix. Although the UI defaults every catalogue permission to granted for Admin, several operational APIs exclude Admin, including routine resident directory, consultations, certificates, and referral workflows. The transfer review API explicitly includes Admin.
\item[MHO.] Municipal access covers resident and referral oversight, health services, official reports, and community analytics. In the certificate service, MHO creation produces a draft for review. MHO can review transfers within its authorized scope. Referral approval and status changes use a separate recipient role gate that excludes MHO. The account-creation catalogue grant is used in resident account workflows. The personnel management API excludes MHO.
\item[PHN.] Municipal scope includes directories and monitoring, clinical processing, follow-ups, referral authoring and assigned recipient decisions, and report submission. The certificate service issues eligible PHN-created certificates at creation. Delegated personnel management allows PHN to view and provision Health Supervisor and RHU Personnel accounts in its municipality, edit permitted profile fields, and deactivate or reactivate eligible accounts. The historical personnel review queue also covers these two roles.
\item[Health Supervisor.] Barangay scope covers resident verification, household verification, consultation, triage, follow-up, health services, and reports. Referral decisions require recipient authorization. RHU Consultation Station access requires an explicit Consultation assignment. Delegated personnel management covers BHW accounts in the assigned barangay, with the same permitted profile and status operations as PHN. The historical approval mapping covers BHW and Resident, with resident registration also handled through its separate verification workflow.
\item[RHU Personnel.] Triage and Consultation workflows are gated by assigned station and facility. Municipal certificate reads use the request grant, while initiating a certificate request requires a Triage assignment. The request is filed for review. Certificate editing and decisions are reserved for PHN and MHO through a separate role gate. Referral accept, reject, and status updates are limited to assigned recipients at the current facility or station. The health-services catalogue API excludes RHU Personnel despite the view grant. The submitted-report API admits RHU Personnel as a scoped recipient and checks \texttt{reports.view} when the role is configured, although that permission is off by default.
\item[BHW.] Assigned-barangay duties cover household profiling and limited resident intake, including identity lookup, new record encoding, permitted edits, and submission for PHN processing. The clinical resident directory and official submitted-report authoring APIs exclude BHW. The default report-generation grant therefore does not authorize official report submission. BHW has separate role-based FHSIS M1 data-entry and barangay appointment-management functions.
\item[Resident.] Access is ownership-scoped. The catalogue limits its grants to self-service viewing. Dedicated resident routes also support booking or cancelling an eligible appointment, responding to a proposed schedule, and submitting or cancelling an eligible residency transfer. Pending verification triggers limited access within the same Resident role.
\end{description}

\noindent\textit{Account management distinction.} The PHN and Health Supervisor delegated account endpoints authorize operations through target-role and assignment rules in the service and database, so their scoped account listing and limited edits can be available even when the corresponding catalogue grants are off. Delegates can change eligible account status through the scoped update workflow. The separate lifecycle, access-reset, deletion, and permission-settings endpoints require Admin.

\par\medskip
\noindent\textit{Source.} KALUSAGAP, branch \texttt{v37}, commit \texttt{9fa1cf3}. Permission definitions are in \texttt{frontend/src/lib/permissions.js}. Backend scope is defined in the API routes and services and in the Supabase migrations.
\par
\endgroup
